(1) 对平面行波，压力与同方向的体速度满足 \(p=Zv\)，其中特征声阻抗 \(Z_j=\rho_jc_j\)。左侧入射压力振幅 1、反射压力 r；对应速度分别为 \(1/Z_1\) 和 \(-r/Z_1\)，右侧透射压力 t 对应速度 \(t/Z_2\)。

(2) 界面压强连续 \(1+r=t\)，体速度连续 \((1-r)/Z_1=t/Z_2\)。解得
\[
r=\frac{Z_2-Z_1}{Z_2+Z_1},\qquad
t=\frac{2Z_2}{Z_1+Z_2}.
\]
注意这些是声压幅度系数，不是能流比；当 Z2>Z1 时压力反射同相。

(3) 一维声波平均能流为 \(\overline P=|p_0|^2/(2Z)\)，所以
\[
\mathcal R=|r|^2=\frac{(Z_2-Z_1)^2}{(Z_1+Z_2)^2},
\quad
\mathcal T=\frac{Z_1}{Z_2}|t|^2=\frac{4Z_1Z_2}{(Z_1+Z_2)^2}.
\]
相加得到 \(\mathcal R+\mathcal T=1\)，无损界面只重新分配方向能流。

(4) 当 Z1=Z2，r=0、\(\mathcal T=1\)，界面在正入射时完全无反射。当 \(Z_2\to\infty\)，r→+1、t→2，但实际穿透功率 \(\mathcal T\to0\)，即刚性高阻端压力加倍、法向速度趋零。
